% CV en español. Compilar con:  pdflatex cv-es.tex
% Copiar el PDF resultante a files/CV_Esteban_Hernandez_Urrutia.pdf
\documentclass[10pt]{article}
\usepackage{cv-style}
\usepackage[spanish]{babel}
\newcommand{\sd}{2.\textsuperscript{do}~sem.}
\newcommand{\pr}{1.\textsuperscript{er}~sem.}

\begin{document}
\cvfooter{Noviembre 2026}{Esteban Hernández Urrutia \,·\, Currículum Vitae}

\cvheader{Esteban}{Hernández Urrutia}{Doctor en Matemáticas}{esteban.hernandez@usm.cl}{ehernandez001-urr.github.io}

\cvsection{Formación académica}
\cventry{Universidad Técnica Federico Santa María y Université Grenoble Alpes}{Valparaíso -- Grenoble}
  {Doctorado en Matemáticas y Doctorat en Automatique (doble grado)}{2017 -- 2021}
  {\item Directores: Eduardo Cerpa (Chile) y Christophe Prieur (Francia)
   \item Tesis: Controlabilidad en la frontera, estabilización y problemas de seguimiento para sistemas parabólicos
   \item Área: Controlabilidad y estabilización de EDP de tipo parabólico}
\cventry{Universidad Técnica Federico Santa María}{Valparaíso, Chile}
  {Magíster en Ciencias Matemáticas}{2014 -- 2017}
  {\item Director: Eduardo Cerpa}
\cventry{Universidad Técnica Federico Santa María}{Valparaíso, Chile}
  {Ingeniería Civil Matemática}{2007 -- 2015}
  {\item Director: Eduardo Cerpa}

\cvsection{Experiencia en investigación}
\cvitem{Ago 2021 -- Dic 2022}{\textbf{Núcleo Milenio de Control Aplicado y Problemas Inversos, Chile}, asistente de investigación, enfocado en control y modelamiento aplicados a neurociencia}
\cvitem{Nov 2019 -- Sep 2020}{\textbf{Gipsa Lab, Université Grenoble Alpes, Francia}, estadía larga de investigación, parte del convenio del programa de doble grado}
\cvitem{Jul 2018}{\textbf{Gipsa Lab, Université Grenoble Alpes, Francia}, estadía de investigación de un mes en el departamento de control automático, bajo la supervisión de Christophe Prieur}
\cvitem{Jul 2017}{\textbf{Gipsa Lab, Université Grenoble Alpes, Francia}, estadía de investigación de un mes en el departamento de control automático, bajo la supervisión de Christophe Prieur}
\cvitem{2011 -- 2013}{\textbf{Universidad Técnica Federico Santa María}, práctica de investigación en control de ecuaciones diferenciales ordinarias, bajo la supervisión de Eduardo Cerpa}

\cvsection{Publicaciones}
\cvsubsection{Publicadas}
% Publications shared by both CVs (titles and journals stay in English)
\cvref{Cerpa, E., Courdurier, M., Hernández, E., Medina, L. E., Paduro, E. 2026. \textit{Approximation and stability results for the parabolic FitzHugh--Nagumo system with combined rapidly oscillating sources}. Discrete and Continuous Dynamical Systems, 50, 69--102. \url{https://doi.org/10.3934/dcds.2026001}}
\cvref{Guzmán, P., Hernández, E. 2025. \textit{Rapid stabilization of an unstable heat equation with disturbance at the flux boundary condition}. Systems \& Control Letters, 196, 105973. \url{https://doi.org/10.1016/j.sysconle.2024.105973}}
\cvref{Guzmán, P., Hernández, E. 2023. \textit{Stabilization of the heat equation with disturbance at the flux boundary condition}. Mathematical Methods in the Applied Sciences, 46(17), 18035--18043. \url{https://doi.org/10.1002/mma.9544}}
\cvref{Cerpa, E., Courdurier, M., Hernández, E., Medina, L. E., Paduro, E. 2023. \textit{A partially averaged system to model neuron responses to interferential current stimulation}. Journal of Mathematical Biology, 86(1), 8. \url{https://doi.org/10.1007/s00285-022-01839-8}}
\cvref{Hernández, E., Prieur, C., Cerpa, E. 2022. \textit{A tracking problem for the state of charge in an electrochemical Li-ion battery model}. Mathematical Control and Related Fields, 12(3), 709--732. \url{https://doi.org/10.3934/mcrf.2021041}}


\cvsection{Formación docente}
\cvrow{Ago -- Sep 2026}{\textbf{Instrumentos de Evaluación por Competencias}, curso aprobado de 28 horas, Dirección de Pregrado}{USM, Valparaíso}

\cvsection{Becas}
\cvitem{2019 -- 2020}{\textbf{PIIC}, Universidad Técnica Federico Santa María}
\cvitem{2017 -- 2021}{\textbf{Beca de Doctorado en Matemáticas}, ANID}
\cvitem{2014 -- 2015}{\textbf{Beca de Magíster en Matemáticas}, ANID}

\cvsection{Charlas}
\cvref{Abril 2024. \textit{Boundary Control Method, una introducción}. Seminario de Control y Problemas Inversos, USM, Valparaíso, Chile.}
\cvref{Julio 2021. \textit{Boundary null controllability of a parabolic--elliptic system}. Webinar de Jovens Pesquisadores de EDP e Matemática Aplicada, Brasil. \url{http://edpemaplicada.uff.br/wpjedpma/}}
\cvref{Abril 2021. \textit{Rapid stabilization of an unstable heat equation with disturbance at the flux boundary condition}. Charla Junior, Webinar: Control in Times of Crisis. \url{http://ctcseminar.mat.utfsm.cl/}}
\cvref{Abril 2021. \textit{Rapid stabilization of an unstable heat equation with disturbance at the flux boundary condition}. Sesión paralela, Jornada de Matemática de la Zona Sur (en línea), Chile. \url{https://www.jmzs.cl/2020/}}
\cvref{Agosto 2020. \textit{A tracking problem for the state of charge in an electrochemical Li-ion battery model}. Charlas de estudiantes, Congreso de Ingeniería PUC--UTFSM (en línea), Chile.}
\cvref{Diciembre 2018. \textit{A tracking problem for the state of charge in an electrochemical Li-ion battery model}. Charlas de estudiantes, ICoPS, Valparaíso, Chile. \url{http://amercado.mat.utfsm.cl/ICOPS2018/}}
\cvref{Agosto 2017. \textit{Adaptive control for parameter identification in parabolic models of Li-ion batteries}. Sesión paralela, COMCA, Arica, Chile.}
\cvref{Noviembre 2016. \textit{Adaptive control for parameter identification in parabolic models of Li-ion batteries}. ENIM, Valparaíso, Chile. \url{http://www.cem.mat.utfsm.cl/enim/actividades/}}
\cvref{Julio 2015. \textit{Control adaptativo de bioprocesos asociados a modelos de microalgas}. Sesión paralela, COMCA, Iquique, Chile.}

\cvsection{Experiencia docente}
\cvrow{\sd{} 2026}{\textbf{MAT052, MAT00225, MATE26}, Profesor}{USM, Valparaíso}
\cvrow{\pr{} 2026}{\textbf{MAT021, MAT051, MAT73425}, Profesor}{USM, Valparaíso}
\cvrow{\sd{} 2025}{\textbf{MAT024, MAT023}, Profesor}{USM, Valparaíso}
\cvrow{\pr{} 2025}{\textbf{MAT012, MAT023}, Profesor}{USM, Valparaíso}
\cvrow{\sd{} 2024}{\textbf{MAT060, MAT070}, Profesor}{USM, Valparaíso}
\cvrow{\sd{} 2024}{\textbf{MAT012, MAT023 ICIPEV}, Profesor}{USM, Valparaíso}
\cvrow{\pr{} 2024}{\textbf{MAT021, MAT023, MAT024}, Profesor}{USM, Valparaíso}
\cvrow{\sd{} 2023}{\textbf{MAT021, MAT022, MATE11, MAT024}, Profesor}{USM, Valparaíso}
\cvrow{\pr{} 2023}{\textbf{MAT021, MAT022, MAT023}, Profesor}{USM, Valparaíso}
\cvrow{\sd{} 2022}{\textbf{MAT021, MAT022, MATE11}, Profesor}{USM, Valparaíso}
\cvrow{\pr{} 2022}{\textbf{MAT021, MAT022, MAT070}, Profesor}{USM, Valparaíso}
\cvrow{\sd{} 2021}{\textbf{MAT021, MAT022}, Profesor}{USM, Valparaíso}
\cvrow{\pr{} 2021}{\textbf{MAT021, MAT022}, Profesor}{USM, Valparaíso}
\cvrow{\sd{} 2018}{\textbf{MAT021, MAT024}, Ayudante}{USM, Valparaíso}
\cvrow{\sd{} 2016}{\textbf{MAT024}, Ayudante}{USM, Valparaíso}
\cvrow{2010 -- 2012}{\textbf{MAT021, MAT022, Probabilidad y Estadística para Ingenieros, Series de Tiempo}, Ayudante}{USM, Valparaíso}

\cvsection{Extensión y desarrollo profesional}
\cvsubsection{Servicio y extensión}
\cvrow{2025}{\textbf{Programa de Inclusión USM, Departamento de Matemática}, Coordinador}{USM, Valparaíso}
\cvrow{Primavera 2024}{\textbf{Programa de Inclusión USM, Departamento de Matemática}, Coordinador}{USM, Valparaíso}
\cvrow{Ene 2023}{\textbf{Escuela de Verano para estudiantes de enseñanza media}, Docente (\url{https://matematica.usm.cl/edmat-2023/})}{USM, Valparaíso}

\end{document}
